% Discussion for item 2.
% scripts/make_latex.py will NEVER overwrite this file -- edit it freely.

\paragraph{2.1 --- the child dies first.}
The child exited immediately, yet \texttt{ps} still listed it as
\texttt{[2\_1\_zombie] <defunct>}, state \texttt{Z}, with PPID 439 --- the
parent's PID. A zombie has already given up its memory, its open files and its
CPU time; all that remains is its entry in the process table, holding its exit
status until the parent collects it.

The cause is the missing \texttt{wait()}. The kernel cannot discard the entry
before then, or the parent could never learn how its child ended. Each zombie is
tiny, but the process table is finite, so a program that forks repeatedly and
never reaps eventually cannot create processes at all --- the same limit item 3
runs into.

\paragraph{2.2 --- the parent dies first.}
The child reported PPID 455 while its parent lived and PPID 301 afterwards: the
kernel re-parented it, because every process must have a parent to report to.

The adopter was \texttt{Relay(303)}, not PID~1 --- even though PID~1 here really
is \texttt{systemd}. A process can mark itself a \emph{child subreaper}, and
orphans are then handed to the nearest such ancestor instead of travelling all
the way up; WSL runs one \texttt{Relay} per session and it claims that role. The
mechanism is the one the lecture describes; only the PID differs.

No \texttt{<defunct>} entry appeared here, because an orphan is not a zombie: it
is alive and merely has a new parent. When it finally exits, that parent reaps it
at once --- which is what a subreaper is for.

\paragraph{Comparing the two.}
Only 2.1 leaks anything. An orphan costs nothing; a zombie persists for as long
as its negligent parent lives. The problem in 2.1 is not that the child died, but
that nobody was listening when it did.
